% ---
% titre: Priorité des opération
% theme: NO
% chapitre: Nombres naturels et décimaux
% sous_chapitre: Priorité des opérations
% niveau: [11e]
% type: EVAL
% tags: [c-priorite-operations,c-puissance,c-racine,p-question-TS]
% difficulte: 1
% corrige: true
% ---

\question[11]
Calcule le résultat des opérations suivantes \textbf{en écrivant les étapes intermédiaires.}

\vspace{30pt}

\begin{parts}
	\part
		$(9 - 4) \cdot 5 + 7 =$
		
	\vspace{10pt}{\footnotesize\raggedleft Espace pour ta démarche et tes calculs\par}
	\begin{solutionorgrid}[100pt]
	
		\vspace{10pt}
		$5 \textbf{\,(0.5pt)} \cdot 5 + 7 = 32 \textbf{\,(0.5pt)}$
	\end{solutionorgrid}	
	
	\vfill	
	\part
		$37 - 72 \div 3^2 + 15 =$
		
	\vspace{10pt}{\footnotesize\raggedleft Espace pour ta démarche et tes calculs\par}
	\begin{solutionorgrid}[100pt]
	
		\vspace{10pt}
		$37 - 72 \div 9 \textbf{\,(0.5pt)} + 15 = 37 - 8 \textbf{\,(0.5pt)} + 15 = 44 \textbf{\,(0.5pt)}$
	\end{solutionorgrid}
	
	\vfill	
	\part
		$(30 - 12) \cdot 3 =$
		
	\vspace{10pt}{\footnotesize\raggedleft Espace pour ta démarche et tes calculs\par}
	\begin{solutionorgrid}[100pt]
	
		\vspace{10pt}
		$18  \textbf{\,(0.5pt)} \cdot 3 = 54 \textbf{\,(0.5pt)}$
	\end{solutionorgrid}	
	
	\vfill	
	\part
		$3 \cdot (10 - 18\div3)^2=$
		
	\vspace{10pt}{\footnotesize\raggedleft Espace pour ta démarche et tes calculs\par}
	\begin{solutionorgrid}[100pt]
	
		\vspace{10pt}
		$3 \cdot (10 - 6 \textbf{\,(0.5pt)}) ^2 = 3 \cdot 4^2 \textbf{\,(0.5pt)} = 3 \cdot 16 \textbf{\,(0.5pt)} = 48 \textbf{\,(0.5pt)}$
	\end{solutionorgrid}
		
	\newpage
	\part
		$\sqrt[3]{125} =$
		
	\vspace{10pt}{\footnotesize\raggedleft Espace pour ta démarche et tes calculs\par}
	\begin{solutionorgrid}[100pt]
	
		\vspace{10pt}
		$5 \textbf{\,(1pt)}$
	\end{solutionorgrid}	
		
	\vfill
	\part
		$15 + (44 \div 4) \cdot 3 - 13 =$
		
	\vspace{10pt}{\footnotesize\raggedleft Espace pour ta démarche et tes calculs\par}
	\begin{solutionorgrid}[100pt]
	
		\vspace{10pt}
		$15 + 11 \textbf{\,(0.5pt)} \cdot 3 - 13 = 15 + 33 \textbf{\,(0.5pt)} -13 = 35 \textbf{\,(0.5pt)}$
	\end{solutionorgrid}	
	
	\vfill	
	\part
		$\sqrt{81} =$
		
	\vspace{10pt}{\footnotesize\raggedleft Espace pour ta démarche et tes calculs\par}
	\begin{solutionorgrid}[100pt]
	
		\vspace{10pt}
		$9 \textbf{\,(1pt)}$
	\end{solutionorgrid}			

	\vfill
	\part
		$89 - 3 \cdot \sqrt{4 \cdot 16} =$
		
	\vspace{10pt}{\footnotesize\raggedleft Espace pour ta démarche et tes calculs\par}
	\begin{solutionorgrid}[100pt]
	
		\vspace{10pt}
		$89 - 3 \cdot \sqrt{64 \textbf{\,(0.5pt)}} = 89 - 3 \cdot 8 \textbf{\,(0.5pt)} = 89 - 24\textbf{\,(0.5pt)} = 65 \textbf{\,(0.5pt)}$
	\end{solutionorgrid}					
	\vfill
\end{parts}
